%%%%%%%%%%%%%%%%%%%%%%%%%%%%%%%%%%%%%%%%%%%%%%%%%%%%%%%%%%%%%%%%%%%%%%
% Lớp 10 | Thống kê | Các số đặc trưng đo xu thế trung tâm của mẫu số liệu
% Dạng 0D6?3-5: Mốt
% Số câu: 21  (sinh tự động bởi tach_cau.py; từ nay thầy sửa trực tiếp file này)
%%%%%%%%%%%%%%%%%%%%%%%%%%%%%%%%%%%%%%%%%%%%%%%%%%%%%%%%%%%%%%%%%%%%%%

%%%%%%%%%%%%%%%% TRẮC NGHIỆM 4 PHƯƠNG ÁN %%%%%%%%%%%%%%%%

% nguồn: MaHoaNH_TracNghiem_Lop10
\begin{ex}%[0D6N3-5]%[ma:3ca89c11]
	Cho mẫu số liêu thống kê $\{6; 5; 5; 2; 9; 10; 8\}$. Mốt của mẫu số liêu trên bằng  bao nhiêu?
	 \choice
		{\True  $5$}
		 {$10$}
		 {$2$}
		{$6$} 
	\loigiai{ 
		Mốt của  mẫu số liêu trên là $M_0=5$.
	} 
\end{ex}

% nguồn: MaHoaNH_TracNghiem_Lop10
\begin{ex}%[0D6N3-5]%[ma:985861e9]
	Mốt của một bảng phân  bố tần số là
	 \choice
		{Tần số lớn nhất trong bảng phân  bố tần số}
		 {\True Giá trị có tần số lớn nhất trong bảng phân bố tần số}
		{Giá  trị có tần số nhỏ nhất trong bảng phân bố tần số}
		{Tần số nhỏ nhất trong bảng  phân bố tần số}
	 \loigiai{
		Mốt của một bảng phân bố tần số  là giá trị có tần số lớn nhất.
	 }
\end{ex}

% nguồn: MaHoaNH_TracNghiem_Lop10
\begin{ex}%[0D6N3-5]%[ma:c1e1a47a]
	Các giá trị xuất hiện nhiều nhất  trong mẫu số liệu được gọi là
	\choice 
		{\True  Mốt}
		{Số trung  bình}
		 {Số trung vị}
		{Độ lệch chuẩn} 
	 \loigiai{
		Các giá trị xuất hiện nhiều  nhất trong mẫu số liệu được gọi là Mốt.
	} 
\end{ex}

% nguồn: MaHoaNH_TracNghiem_Lop10
\begin{ex}%[0D6N3-5]%[ma:e78f80da]
	Số  áo bán được trong một quý ở cửa hàng bán áo sơ mi nam được thống kê như sau
	\begin{longtable}{|c|c|c|c|c|c|c|c|}
		\hline Cỡ áo & $36$ & $37$ & $38$ & $39$ & $40$ & $41$ & $42$  \\
		\hline Số áo bán được & $13$ & $45$ & $126$ & $125$ &  $150$ & $40$ & $12$\\
		 \hline
	\end{longtable}
	Giá  trị mốt của mẫu số liệu trên bằng
	\choice 
		 {$ 150 $}
		{\True  $ 40 $}
		 {$ 12 $}
		{$ 42  $}
	 \loigiai{
		Vì tần số lớn nhất là $ 150 $ và $ 150 $ tương ứng với cỡ áo $ 40  $ nên mốt của bảng số liệu trên là $ 40 $.
	 }
\end{ex}

% nguồn: MaHoaNH_TracNghiem_Lop10
\begin{ex}%[0D6N3-5]%[ma:138756e0]
	Kết quả thi môn Toán giữa kì 1$1$  của lớp $10A_3$ trường THPT Ba Vì được thống kê như sau
	\begin{center}
		\begin{tabular}{|c|c|c|c|c|c|c|c|}
			\hline Điểm thi &  $5$ & $6$ & $7$ & $8$ & $9$ & $10$ & Cộng\\
			\hline Tần số & $5$ & $7$  & $8$ & $12$ & $8$ & $5$ & $45$\\
			 \hline
		\end{tabular}
	\end{center}
	Giá trị mốt $M_0$ của bảng phân bố tần số  trên bằng
	 \choice
		{$5$} 
		 {$7$}
		{\True  $8$}
		 {$12$}
	 \loigiai{
		Mốt của bảng phân bố tần suất là giá trị có tần số lớn nhất nên ta  có $M_0=8$.
	} 
\end{ex}

% nguồn: MaHoaNH_TracNghiem_Lop10
\begin{ex}%[0D6N3-5]%[ma:388507eb]
	Tiền thưởng (triệu đồng) của cán bộ và nhân viên trong một công  ty được cho ở bảng dưới đây
	\begin{center}
		\begin{tabular}{|c|c|c|c|c|c|c|}
			\hline Tiền thưởng & $12$ & $13$  & $14$ & $15$ & $16$ & Cộng\\
			\hline Tần số & $25$ & $15$  & $11$ & $16$ & $17$ & $84$\\
			 \hline
		\end{tabular}
	\end{center}
	Tính  mốt $M_0$.
	\choice 
		 {$M_0=15$}
		 {\True $M_0=12$}
		 {$M_0=10$}
		{$M_0=16$} 
	\loigiai{ 
		Vì tần số  bằng $25$ là lớn nhất và ứng với số tiền thưởng (triệu đồng) của cán bộ và nhân viên trong một công ty là $12$ nên $M_0=12$.
	} 
\end{ex}

% nguồn: MaHoaNH_TracNghiem_Lop10
\begin{ex}%[0D6N3-5]%[ma:8cf49b1c]
	Thống kê $5$ size áo sơ mi của một cửa hàng thời trang trong tháng $4$ được cho trong  bảng sau
	\begin{center}
		\begin{tabular}{|c|c|c|c|c|c|c|}
			\hline Size áo & 38 & 39 & 40 &  41 & 42 & Cộng \\
			\hline Tần số & 3 & 6 & 15 &  10 & 6 & 40 \\
			 \hline
		\end{tabular}
	\end{center}
	Mốt của  mẫu số liệu trên là
	\choice 
		 {$M_\circ=42$}
		{$M_\circ=38$} 
		{\True $M_\circ=40$} 
		 {$M_\circ=15$}
	\loigiai{ 
		Mốt của mẫu số  liệu trên là $M_\circ=40$.
	 }
\end{ex}

% nguồn: MaHoaNH_TracNghiem_Lop10
\begin{ex}%[0D6N3-5]%[ma:5e754439]
	Cho bảng số liệu nghi lại điểm của $40$ học sinh trong bài kiểm  tra $1$ tiết môn Toán.
	\begin{center}
		\begin{tabular}{|c|c|c|c|c|c|c|c|c|c|}
			\hline 
			Điểm & $3$ &  $4$ & $5$ & $6$ & $7$ & $8$ & $9$ & $10$ & Cộng\\
			 \hline
			Số học sinh & $2$ & $3$ & $7$ & 18 & $3$ & $2$ & $4$  & $1$ & $40$\\
			 \hline
		\end{tabular}
	\end{center}
	 Mốt của dấu hiệu
	 \choice
		 {$M_0=10$}
		 {$M_0=18$}
		{\True  $M_0=6$}
		 {$M_0=40$}
	 \loigiai{
		Điểm $6$ xuất hiện nhiều  nhất nên $M_0=6$.
	 }
\end{ex}

% nguồn: MaHoaNH_TracNghiem_Lop10
\begin{ex}%[0D6N3-5]%[ma:288301fa]
	Một  cửa hàng trà sữa vừa khai trương, thống kê lượng khách tới quán trong $7$ ngày đầu và thu được mẫu số liệu sau
	\begin{center}
		\begin{tabular}{|c|c|c|c|c|c|c|}
			 \hline
			Ngày $1$ & Ngày $2$ & Ngày $3$ & Ngày $4$ & Ngày $5$ & Ngày $6$ &  Ngày $7$\\
			 \hline
			$575$ & $454$ & $400$ & $325$ & $351$  & $333$ & $412$\\
			\hline 
		\end{tabular}
	\end{center}
	Mệnh đề  nào sau đây là đúng?
	 \choice
		{Số trung  vị là $263$}
		{\True Số trung bình làm tròn đến hàng  phần trăm là $407{,}14$}
		 {Dấu hiệu điều tra ở đây là doanh thu của quán trà sữa}
		{Ngày $2$ là mốt của  mẫu số liệu này}
	 \loigiai{
		Ta có $\overline  {x}=\dfrac{575+454+400+325+351+333+412}{7}=407{,}14$.
	} 
\end{ex}

% nguồn: MaHoaNH_TracNghiem_Lop10
\begin{ex}%[0D6N3-5]%[ma:145f44f7]
	Số áo bán được trong một quý ở cửa hàng  bán áo sơ mi nữ được thống kê như sau:
	\begin{center}
		\begin{tabular}{|l|c|c|c|c|c|c|c|}
			\hline Cỡ áo & 36 &  37 & 38 & 39 & 40 & 41 & 42 \\
			 \hline
			 Tần số
			 (Số áo bán được)
			& 13 & 45  & 126 & 125 & 110 & 40 & 12 \\
			\hline 
		\end{tabular}
	\end{center}
	Giá trị mốt  của bảng phân bố tần số trên bằng
	 \choice
		{\True $38$} 
		{$126$} 
		 {$42$}
		 {$12$}
	 \loigiai{
	Vì giá trị $x_3=38$ có tần số $n_3=126$ lớn nhất  nên mốt của bảng số liệu là $M_0=38$.}
\end{ex}

% nguồn: MaHoaNH_TracNghiem_Lop10
\begin{ex}%[0D6N3-5]%[ma:64648e59]
	Cho bảng phân bố tần số tiền thưởng (triệu đồng) cho cán bộ và nhân  viên trong một công ty
	\begin{center}
		\begin{tabular}{|c|c|c|c|c|c|c|}
			\hline 
			Tiền thưởng & $2$ & $3$ & $4$ & $5$  & $6$ & Cộng\\
			 \hline
			Tần  số & $5$ & $15$ & $10$ & $6$ & $7$ & $43$\\
			 \hline
		\end{tabular}
	\end{center}
	Mốt của bảng phân bố  tần số đã cho là
	\choice 
		 {\True $3$ triệu đồng}
		{$2$ triệu đồng} 
		{$6$ triệu đồng} 
		 {$5$ triệu đồng}
	 \loigiai{
		Ta có $M_0=3$. 
	} 
\end{ex}

% nguồn: MaHoaNH_TracNghiem_Lop10
\begin{ex}%[0D6N3-5]%[ma:be94accc]
	Điều tra tiền lương một tháng của $100$ người lao động trên địa bàn một xã ta  có bảng phân bố tần số sau:
	\begin{center}
		\begin{tabular}{|c|c|c|c|c|c|c|}
			 \hline
			\begin{tabular}{c}
				 Tiền lương \\
				(VND) 
			\end{tabular}
			& 6 000 000 & 7 000 000 & 8 500 000 & 9 500 000 &  10 000 000 & 10 500 000 \\
			\hline Tần số & 13 & 37 &  25 & 11 & 6 & 8 \\
			\hline 
		\end{tabular}
	\end{center}
	Tìm mốt  của bảng phân bố tần số trên.
	 \choice
		 {$6\ 000\ 000$}
		{\True $7\  000\ 000$}
		 {$8\ 500\ 000$}
		{$9\  500\ 000$}
	\loigiai{ 
	Ta có giá  trị $7\ 000\ 000$ có tần số lớn nhất nên là mốt của bảng phân bố tần số trên.}
\end{ex}

% nguồn: MaHoaNH_TracNghiem_Lop10
\begin{ex}%[0D6N3-5]%[ma:191a4a4c]
	Cho bảng phân  bố tần số như sau
	\begin{center}
		\begin{tabular}{|c|c|c|c|c|c|c|c|c|}
			\hline Giá trị & $x_1$ & $x_2$ & $x_3$ & $x_4$ &  $x_5$ & $x_6$ & $x_7$ & $x_8$\\
			\hline Tần số & $15$ & $9 n-1$ & $12$ & $n^2+7$ & $14$ &  $10$ & $9 n-20$ & $17$\\
			\hline 
		\end{tabular}
	\end{center}
	 Tìm $n$ để $M_{01}=x_2$; $M_{02}=x_4$ là hai mốt của bảng số liệu trên.
	 \choice
		 {$n=1$; $n=8$}
		{\True  $n=8$}
		 {$n=1$}
		{$n=9$} 
	\loigiai{ 
		Ta có $M_{01}=x_2$; $M_{02}=x_4$ là hai mốt của bảng phân bố tần  số nên
		\allowdisplaybreaks 
		\begin{eqnarray*}
			\heva{&n^2+7=9n-1\\&9n-1 > 17} \Leftrightarrow \heva{&n^2-9n+8=0\\&n > 2}  \Leftrightarrow \heva{&\hoac{&n=1\,(\text{loại})\\&n=8\,(\text{thỏa mãn})}\\
			&n > 2} \Rightarrow  n=8.
		\end{eqnarray*}
	} 
\end{ex}


%%%%%%%%%%%%%%%% ĐÚNG/SAI %%%%%%%%%%%%%%%%

% nguồn: MaHoaNH_DungSai_Lop10
\begin{ex}%[0D6N3-5]%[ma:1c2567ea]
	Cho mẫu số liệu  sau: $4$; $5$; $6$; $7$; $8$; $4$; $9$; $4$; $3$. Xét tính \textbf{đúng}, \textbf{sai} của mệnh đề
	\choiceTF 
		{\True Số trung bình  $\overline {x}=5{,}5$}
		{Mốt $M_0=3$} 
		{Trung vị là  $M_e=4$}
		{\True Tứ phân vị thứ ba là  $Q_3=7$}
	 \loigiai{
		Sắp xếp mẫu: $3$; $4$; $4$; $4$; $5$; $5$; $6$; $7$; $8$;  $9$.\\
		Số trung  bình $\overline {x}=\dfrac{55}{10}=5{,}5$.\\
		Kích thước mẫu là $10$ (chẵn) nên trung vị là $M_e=\dfrac{1}{2}(5+5)=5$.\\ 
		 Tứ phân vị thứ nhất là $Q_1=4$.\\
		Tứ  phân vị thứ hai là $Q_2=M_e=5$.\\
		 Tứ phân vị thứ ba là $Q_3=7$.\\
		 Mốt $M_0=4$.\\
		Do  đó
		\begin{itemchoice}
			\itemch  Đúng.
			\itemch Sai. 
			\itemch Sai. 
			 \itemch Đúng.
		\end{itemchoice}
	 }
\end{ex}

% nguồn: MaHoaNH_DungSai_Lop10
\begin{ex}%[0D6H3-5]%[ma:330acd96]
	Cho bảng số liệu điểm  kiểm tra môn Toán của $20$ học sinh.
	\begin{center}
		\begin{tabular}{|c|c|c|c|c|c|c|c|c|}
			\hline Điểm & 4 & 5 & 6 &  7 & 8 & 9 & 10 & Cộng \\
			\hline Số học sinh & 1 & 2 & 3 & 4 & 5 & 4 & 1 & 20 \\ 
			 \hline
		\end{tabular}
	\end{center}
	Các mệnh  đề sau đúng hay sai?
	 \choiceTF
		{Số  trung bình $\overline{x}=7{,}5$}
		{\True Tứ phân  vị thứ hai là $Q_2=7{,}5$}
		{\True Tứ  phân vị thứ nhất là $Q_1=6$}
		 {Mốt là $M_0=10$}
	\loigiai{ 
		\begin{itemchoice}
			\itemch Sai. Số trung bình là $\overline{x}=\dfrac{4+2\cdot 5+6\cdot 3+7\cdot 4+8\cdot  5+9\cdot 4+10}{20}=7{,}3$.
			 \itemch Đúng. Vì cỡ mẫu là $n=20$, là số chẵn, nên giá trị của tứ phân vị thứ hai là $Q_2=\dfrac{7+8}{2}=7{,}5$.
			\itemch Đúng. Tứ phân vị thứ nhất là  $Q_1=\dfrac{6+6}{2}=6$.
			\itemch Sai. Vì $8$ có tần số lớn nhất  nên $M_0=8$.
		\end{itemchoice}
	 }
\end{ex}

% nguồn: MaHoaNH_DungSai_Lop10
\begin{ex}%[0D6H3-5]%[ma:39ff0f5b]
	Điểm điều tra  về chất lượng sản phẩm mới (thang điểm $100$) như sau:
	\begin{center}
		\begin{tabular}{lllllllllllll}72 & 68 & 39 & 41 & 54 & 61 & 72 & 75 & 80 & 61 & 50 & 65 & 78
		\end{tabular}
	\end{center}
	Các mệnh đề sau đúng  hay sai?
	\choiceTF 
		{$n=12$} 
		{\True Mốt của mẫu  số liệu là $61$}
		{Tứ phân vị thứ hai là số  thứ $6$}
		{Tổng số trung vị, tứ phân vị thứ nhất và tứ phân vị thứ  ba là $190$}
	\loigiai{\begin{itemchoice}
	\itemch Sai. Mẫu số liệu  có $13$ số liệu nên $n=13$.
	\itemch Đúng. Vì điểm  $61$ xuất hiện $2$ lần với tần suất lớn nhất nên mốt của mẫu số liệu là $61$.
	\itemch Sai. Vì $n=13$ là số lẻ nên số trung vị là  số đứng ở vị trí thứ $\dfrac{13+1}{2}=7$.
	\itemch Sai. Sắp xếp mẫu số liệu theo  thứ tự tăng dần ta có
	\begin{center}
		\begin{tabular}{lllllllllllll}
			39 & 41 & 50 & 54 & 61 &  61 & 65 & 68 & 72 & 72 & 75 & 78 & 80.
		\end{tabular}
	\end{center}
	 Số trung vị là số đứng thứ $7$ nên $Q_2=65$.\\
	Xét nửa mẫu số liệu bên trái: $\begin{array}{llllll}39 & 41 & 50 & 54 & 61 & 61.
\end{array}$\\
Tứ phân vị thứ nhất chính là trung vị của mẫu số liệu  này.\\ Như thế $Q_1=\dfrac{50+54}{2}=52$.\\
Xét nửa mẫu số liệu bên phải: $\begin{array}{llllll}68 & 72 & 72 & 75 & 78 & 80.
\end{array}$\\
Tứ phân vị thứ ba chính là trung vị nửa  mẫu số liệu này. Do đó $Q_3=\dfrac{72+75}{2}=73{,}5$.\\
Tổng số trung vị, tứ phân vị thứ nhất và tứ phân vị thứ ba là $65+52+73{,}5=190{,}5$. 
\end{itemchoice}
 }
\end{ex}


%%%%%%%%%%%%%%%% TRẢ LỜI NGẮN %%%%%%%%%%%%%%%%

% nguồn: MaHoaNH_TraLoiNgan_Lop10
\begin{bt}%[0D6N3-5]%[ma:3fd1394a]
	Có $100$ học sinh tham gia  kì thi khảo sát chất lượng môn Toán. Điểm khảo sát được tính theo thang điểm $10$ và thống kê như sau:
	\begin{center}
		\begin{tabular}{|l|c|c|c|c|c|c|c|c|c|c|}
			\hline Điểm & 1 & 2 & 3 & 4 &  5 & 6 & 7 & 8 & 9 & 10 \\
			\hline Số & 11 & 8 &  41 & 56 & 29 & 34 & 61 & 50 & 91 & 19 \\
			\hline 
		\end{tabular}
	\end{center}
	Mốt  của mẫu số liệu trên bằng bao nhiêu?
	\shortans{9} 
	 \loigiai{
		Số học sinh đạt điểm  $9$ là nhiều nhất ($91$ học sinh) nên mốt của mẫu số liệu là $9$.
	} 
\end{bt}

% nguồn: MaHoaNH_TraLoiNgan_Lop10
\begin{bt}%[0D6H3-5]%[ma:2c11a084]
	Điểm kiểm tra môn Toán cuối năm của một nhóm gồm $9$ học sinh lớp $10$ lần lượt là $3$; $4$; $5$; $6$; $7$; $8$; $8$; $9$; $9$. Tứ phân vị thứ nhất của mẫu số liệu trên bằng bao nhiêu? (làm  tròn kết quả đến hàng đơn vị)
	 \shortans{4}
	\loigiai{Số điểm của $9$ học  sinh đã được sắp xếp theo thứ tự không giảm
		$$3;\ 4;\  5;\ 6;\ 7;\ 8;\ 8;\ 9;\ 9.$$
		Vì cỡ mẫu  là số lẻ $n=9$ nên tứ phân vị thứ hai là $Q_2=7$.\\
	Tứ  phân vị thứ nhất là trung vị của mẫu: $3$; $4$; $5$, suy ra $Q_1=4$.}
\end{bt}

% nguồn: MaHoaNH_TraLoiNgan_Lop10
\begin{bt}%[0D6H3-5]%[ma:bd515969]
	Thời gian (đơn vị giờ) dành cho hoạt động thể thao trong  tuần của một số học sinh được thống kê như sau:
	\begin{center}
		\begin{tabular}{|l|l|l|l|l|l|l|l|l|l|}
			\hline 1 & 0 & 3 & 2 & 1 & 1 & 4 & 1,5 & 2 & 4  \\
			 \hline
		\end{tabular}
	\end{center}
	Tính mốt của mẫu số liệu trên. 
	\shortans{$1$} 
	\loigiai{ 
		Vì số học sinh dành  $1$ giờ mỗi tuần cho hoạt động thể thao là nhiều nhất ($3$ học sinh) nên mốt là $1$.
	} 
\end{bt}

% nguồn: MaHoaNH_TraLoiNgan_Lop10
\begin{bt}%[0D6H3-5]%[ma:ba27c423]
	Số điểm mà $5$ vận động viên bóng rổ ghi được trong một trận đấu: 
	\begin{center}
		\begin{tabular}{ccccc}
			9 & 8  & 15 & 8 & 20.
		\end{tabular}
	\end{center}
	Tính  tứ phân vị thứ nhất của mẫu số liệu.
	 \shortans{$8$}
	 \loigiai{
		 Sắp xếp lại mẫu số liệu
		\begin{center}
			\begin{tabular}{ccccc}
				8 & 8 &  9 & 15 & 20.
			\end{tabular}
		\end{center}
		Tứ phân vị thứ nhất là $Q_1=\dfrac{8+8}{2}=8$. 
	 }
\end{bt}

% nguồn: MaHoaNH_TraLoiNgan_Lop10
\begin{bt}%[0D6V3-5]%[ma:2ab5d073]
	Cho  bảng phân bố tần số như sau
	\begin{longtable}{|c|c|c|c|c|c|c|c|c|}
		\hline 
		Giá trị&$x_1$&$x_2$&$x_3$&$x_4$&$x_5$&$x_6$&$x_7$&$x_8$\\ 
		\hline 
		Tần  số&$15$&$2n-1$&$12$&$n^2-14n+47$&$14$&$10$&$16$&$17$\\
		\hline 
	\end{longtable}
	Tìm $n$ để  $M_0^{(1)}=x_2$; $M_0^{(2)}=x_4$ là hai mốt của bảng số liệu trên.
	\shortans{12} 
	 \loigiai{
		Ta có $M_0^{(1)}$ và $M_0^{(2)}$ là hai mốt của bảng số liệu trên nên $M_0^{(1)}=M_0^{(2)}\Rightarrow  x_2=x_4$\\
		$\Leftrightarrow  \heva{&2n-1>17\\&2n-1=n^2-14n+47}\Leftrightarrow \heva{&n>9\\&\hoac{&n=4\\&n=12}}\Rightarrow n=12$.
	 }
\end{bt}
